\documentclass[11pt]{article}

\usepackage[margin=0.9in]{geometry}
\usepackage{amsmath,amssymb}
\usepackage{booktabs}
\usepackage{enumitem}
\usepackage{listings}
\usepackage{xcolor}
\usepackage{hyperref}

\definecolor{backcolour}{rgb}{0.95,0.95,0.92}
\definecolor{codeblue}{rgb}{0.05,0.15,0.55}
\definecolor{codegreen}{rgb}{0.0,0.45,0.0}

\lstset{
    backgroundcolor=\color{backcolour},
    basicstyle=\ttfamily\small,
    keywordstyle=\color{codeblue}\bfseries,
    commentstyle=\color{codegreen},
    stringstyle=\color{purple},
    breaklines=true,
    frame=single,
    showstringspaces=false,
    tabsize=4
}

\title{CPSC 250 \\ Fibonacci Sequence and Algorithm Complexity}
\author{Edward Brash}
\date{}

\begin{document}

\maketitle

\section{The Fibonacci Sequence}

The Fibonacci sequence is defined by:

\[
0, 1, 1, 2, 3, 5, 8, 13, 21, \dots
\]

The sequence begins with:

\[
f_0 = 0
\]
\[
f_1 = 1
\]

and every later term is the sum of the previous two:

\[
f_n = f_{n-1} + f_{n-2}
\]

\section{Recursive Fibonacci}

The recursive definition naturally suggests a recursive function.

\begin{lstlisting}[language=Python]
def fibonacci(n):

    if n == 0:
        return 0

    elif n == 1:
        return 1

    else:
        return fibonacci(n - 1) + fibonacci(n - 2)
\end{lstlisting}

\section{Testing the Recursive Function}

\[
fibonacci(0) = 0
\]

\[
fibonacci(1) = 1
\]

\[
fibonacci(2) = fibonacci(1) + fibonacci(0) = 1
\]

\[
fibonacci(3) = fibonacci(2) + fibonacci(1) = 2
\]

The recursive implementation works correctly.

\section{The Computational Problem}

Although recursive Fibonacci is elegant, it is extremely inefficient.

Why?

Because the function repeatedly recalculates the same values.

For example:

\[
fibonacci(5)
\]

requires:

\[
fibonacci(4)
\]

and:

\[
fibonacci(3)
\]

But:

\[
fibonacci(4)
\]

also requires another computation of:

\[
fibonacci(3)
\]

and so on.

\section{Function Call Explosion}

The number of recursive calls grows very rapidly.

Approximate call counts:

\begin{center}
\begin{tabular}{cc}
\toprule
$n$ & Number of Calls \\
\midrule
0 & 1 \\
1 & 1 \\
2 & 3 \\
3 & 4 \\
4 & 7 \\
5 & 11 \\
6 & 18 \\
7 & 29 \\
8 & 47 \\
9 & 76 \\
\bottomrule
\end{tabular}
\end{center}

By:

\[
n = 30
\]

the number of recursive calls is enormous.

\section{Complexity of Recursive Fibonacci}

The recursive algorithm has approximately exponential complexity:

\[
O(e^n)
\]

This is extremely bad scaling behavior.

As \(n\) increases, runtime grows explosively.

\section{Non-Recursive Fibonacci}

We can do much better using iteration.

\begin{lstlisting}[language=Python]
def fibonacci(n):

    if n == 0:
        return 0

    elif n == 1:
        return 1

    else:

        t1 = 0
        t2 = 1

        for i in range(2, n + 1):

            t3 = t2 + t1

            t1 = t2
            t2 = t3

        return t3
\end{lstlisting}

\section{Complexity of Iterative Fibonacci}

The iterative version only performs one loop over the data.

Its complexity is:

\[
O(n)
\]

This is dramatically better than exponential growth.

\section{The Golden Ratio}

The ratio of successive Fibonacci numbers approaches a constant:

\[
\phi \approx 1.618033988\ldots
\]

called the golden ratio.

Using:

\[
f_n = f_{n-1} + f_{n-2}
\]

and dividing by \(f_{n-1}\):

\[
\frac{f_n}{f_{n-1}}
=
1
+
\frac{f_{n-2}}{f_{n-1}}
\]

As \(n\) becomes large:

\[
\phi = 1 + \frac{1}{\phi}
\]

This gives:

\[
\phi^2 - \phi - 1 = 0
\]

Using the quadratic formula:

\[
\phi =
\frac{1 + \sqrt{5}}{2}
\]

\section{Closed-Form Fibonacci Formula}

Binet's formula gives a direct expression for Fibonacci numbers:

\[
f_n =
\frac{
\phi^n -
\left(-\frac{1}{\phi}\right)^n
}{
\sqrt{5}
}
\]

This suggests a constant-time Fibonacci algorithm.

\section{Formula-Based Fibonacci}

\begin{lstlisting}[language=Python]
import math

phi = (1.0 + math.sqrt(5.0)) / 2.0

def fibonacci(n):

    return int(
        (phi**n - (-1.0 / phi)**n)
        / math.sqrt(5.0)
    )
\end{lstlisting}

\section{Floating Point Precision Problems}

For large \(n\), floating-point precision errors appear.

Why?

Because numbers like:

\[
\sqrt{5}
\]

and:

\[
\phi
\]

cannot be represented exactly in binary floating-point form.

Eventually the rounding errors accumulate.

\section{Using the Decimal Module}

Python's \texttt{decimal} module allows arbitrary precision arithmetic.

\begin{lstlisting}[language=Python]
import decimal as dec

dec.getcontext().prec = 50
\end{lstlisting}

Now calculations can be performed with much higher precision.

\section{Summary of Fibonacci Approaches}

\begin{center}
\begin{tabular}{ccc}
\toprule
Method & Complexity & Notes \\
\midrule
Recursive & $O(e^n)$ & Elegant but slow \\
Iterative & $O(n)$ & Efficient and practical \\
Formula-Based & $O(1)$ & Fast but precision-sensitive \\
\bottomrule
\end{tabular}
\end{center}

\section{Key Takeaways}

\begin{itemize}
    \item Recursive algorithms can be elegant but inefficient.
    \item Complexity describes how runtime scales with input size.
    \item Exponential complexity becomes unusable quickly.
    \item Iterative solutions are often dramatically faster.
    \item Floating-point precision matters in numerical computing.
    \item Python's \texttt{decimal} module supports high-precision arithmetic.
\end{itemize}

\end{document}
